\chapter{Исследовательская часть}

\section{Модель вычислений}

Для оценки трудоёмкости алгоритмов умножения матриц используется модель вычислений, основанная на подсчёте элементарных операций:
\begin{enumerate}
    \item операции присваивания, сложения, вычитания (включая унарные знаки), побитового сдвига, взятия индекса, инкремента и сравнения ($=$, $+$, $-$, $+=$, $-=$, $\ll$, $\gg$, $[]$, $++$, $--$, $<$, $>$, $<=$, $>=$, $==$, $!=$) имеют единичную трудоёмкость: $1$;
    \item операции умножения, деления и остатка от деления ($*$, $/$, $\%$, $*=$, $/=$, $\%=$) имеют трудоёмкость: $2$;
    \item трудоёмкость условного оператора определяется как сумма трудоёмкости вычисления условия и трудоёмкости выполняемой ветви:
    \begin{equation}
        f_{\text{if}} = f_{\text{усл}} + \begin{cases} 
            f_{\text{ветви}}, & \text{ветвь выполняется}, \\
            0, & \text{ветвь не выполняется};
        \end{cases}
    \end{equation}
    \item трудоёмкость цикла $\texttt{for}$ с числом итераций $K$ вычисляется с учётом инициализации счётчика, проверки условия продолжения цикла на каждой итерации и шага инкремента:
    \begin{equation}
        f_{\text{for}} = f_{\text{иниц}} + f_{\text{усл}} + K \cdot (f_{\text{усл}} + f_{\text{инкр}} + f_{\text{тела}}).
    \end{equation}
    Для заголовка \texttt{for (size\_t i = 0; i < N; ++i)}:
    \begin{equation}
        f_{\text{for}} = 1 + 1 + K \cdot (1 + 1 + f_{\text{тела}}) = 2 + K \cdot (2 + f_{\text{тела}}).
    \end{equation}
    Если условие цикла содержит целочисленное деление, например \texttt{k < N / 2}, то $f_{\text{усл}} = 3$. В этом случае:
    \begin{equation}
        f_{\text{for}} = 1 + 3 + K \cdot (3 + 1 + f_{\text{тела}}) = 4 + K \cdot (4 + f_{\text{тела}});
    \end{equation}
    \item затраты на вызовы функций, динамическое выделение памяти под матрицы и получение их размеров в оценку не включаются.
\end{enumerate}

Обозначение: $H = \lfloor N / 2 \rfloor$. Для квадратных матриц размера $N \times N$ при чётных значениях $N$ $H = N / 2$, а при нечётных --- $H = (N - 1) / 2$.

\section{Оценка трудоёмкости алгоритмов}

\subsection{Стандартный алгоритм умножения матриц}

Тело цикла по $k$ содержит оператор:
\begin{center}
\texttt{C[i][j] += A[i][k] * B[k][j];}
\end{center}
Трудоёмкость тела цикла по $k$: $f_{\text{тела\_}k} = 6 + 2 + 1 = 9$.

Условие цикла \texttt{k < N} ($f_{\text{усл}} = 1$), число итераций равно $N$. Затраты на выполнение цикла по $k$:
\begin{equation}
    f_k = 2 + N \cdot (2 + 9) = 11N + 2.
\end{equation}

Цикл по $j$ выполняется $Q$ раз с условием ($f_{\text{усл}} = 1$):
\begin{equation}
    f_j = 2 + Q \cdot (2 + f_k) = 2 + Q \cdot (11N + 4) = 11NQ + 4Q + 2.
\end{equation}

Цикл по $i$ выполняется $M$ раз с условием ($f_{\text{усл}} = 1$):
\begin{equation}
    f_{\text{станд}}(M, N, Q) = 2 + M \cdot (2 + f_j) = 11MNQ + 4MQ + 4M + 2.
\end{equation}

Для квадратных матриц ($M = N = Q$) трудоёмкость стандартного алгоритма составляет:
\begin{equation}
    f_{\text{станд}}(N) = 11N^3 + 4N^2 + 4N + 2 = O(N^3).
\end{equation}

\subsection{Базовый алгоритм Винограда}

Трудоёмкость алгоритма складывается из четырёх этапов:
\begin{equation}
    f_{\text{вин}} = f_{row} + f_{col} + f_{\text{осн}} + f_{\text{нечёт}}.
\end{equation}

При вычислении вектора $row$ тело внутреннего цикла:
\begin{center}
\texttt{row[i] += A[i][2 * j] * A[i][2 * j + 1];}
\end{center}
Итого $f_{\text{тела\_}row} = 1 + 4 + 1 + 4 + 2 + 1 = 13$. Внутренний цикл имеет условие \texttt{j < N / 2} с делением ($f_{\text{усл}} = 3$) и совершает $H$ итераций:
\begin{equation}
    f_{j\_row} = 4 + H \cdot (4 + 13) = 17H + 4.
\end{equation}
Внешний цикл по $i$ ($M$ итераций, условие \texttt{i < M}, $f_{\text{усл}} = 1$):
\begin{equation}
    f_{row} = 2 + M \cdot (2 + 17H + 4) = 17MH + 6M + 2.
\end{equation}

Структура вычисления вектора $col$ аналогична вектору $row$: внешний цикл по $i$ совершает $Q$ итераций, внутренний --- $H$ итераций с условием \texttt{j < N / 2} ($f_{\text{тела\_}col} = 13$, $f_{\text{усл}} = 3$):
\begin{equation}
    f_{col} = 17QH + 6Q + 2.
\end{equation}

В основной части в цикле по $j$ выполняется инициализация:
\begin{center}
\texttt{C[i][j] = -row[i] - col[j];}
\end{center}
Итого $f_{\text{иниц\_}C} = 7$.

Тело цикла по $k$:
\begin{center}
\texttt{C[i][j] += (A[i][2 * k] + B[2 * k + 1][j]) * (A[i][2 * k + 1] + B[2 * k][j]);}
\end{center}
Итого $f_{\text{тела\_}k} = 2 + 8 + 2 + 8 + 2 + 2 + 1 = 25$.

Цикл по $k$ выполняется $H$ раз с условием \texttt{k < N / 2} ($f_{\text{усл}} = 3$):
\begin{equation}
    f_{k\_\text{осн}} = 4 + H \cdot (4 + 25) = 29H + 4.
\end{equation}
Цикл по $j$ ($Q$ итераций, $f_{\text{усл}} = 1$):
\begin{equation}
    f_{j\_\text{осн}} = 2 + Q \cdot (2 + f_{\text{иниц\_}C} + f_{k\_\text{осн}}) = 2 + Q \cdot (2 + 7 + 29H + 4) = 29QH + 13Q + 2.
\end{equation}
Цикл по $i$ ($M$ итераций, $f_{\text{усл}} = 1$):
\begin{equation}
    f_{\text{осн}} = 2 + M \cdot (2 + 29QH + 13Q + 2) = 29MQH + 13MQ + 4M + 2.
\end{equation}

Затраты на проверку условия \texttt{N \% 2 == 1}: $f_{\text{усл\_нечёт}} = 3$. При чётном $N$ ветвь не выполняется: $f_{\text{нечёт}} = 3$. При нечётном $N$ выполняется цикл с телом \texttt{C[i][j] += A[i][N - 1] * B[N - 1][j];} ($f_{\text{тела}} = 11$, $f_j = 13Q + 2$, $f_i = 13MQ + 4M + 2$):
\begin{equation}
    f_{\text{нечёт}} = 3 + 13MQ + 4M + 2 = 13MQ + 4M + 5.
\end{equation}

Трудоёмкость базового алгоритма Винограда:
\begin{equation}
    f_{\text{вин}} = 29MQH + 17H(M + Q) + 13MQ + 10M + 6Q + 6 + f_{\text{нечёт}}.
\end{equation}

Для квадратных матриц ($M = N = Q$):
\begin{itemize}
    \item при чётном $N$ ($H = N / 2$, $f_{\text{нечёт}} = 3$):
    \begin{align}
        f_{\text{вин\_чёт}}(N) &= 29N^2 \cdot \frac{N}{2} + 17 \cdot \frac{N}{2} \cdot 2N + 13N^2 + 16N + 9 \notag \\
        &= 14{,}5N^3 + 30N^2 + 16N + 9;
    \end{align}
    \item при нечётном $N$ ($H = (N - 1) / 2$, $f_{\text{нечёт}} = 13N^2 + 4N + 5$):
    \begin{align}
        f_{\text{вин\_нечёт}}(N) &= 29N^2 \cdot \frac{N - 1}{2} + 17 \cdot \frac{N - 1}{2} \cdot 2N + 26N^2 + 20N + 11 \notag \\
        &= 14{,}5N^3 + 28{,}5N^2 + 3N + 11.
    \end{align}
\end{itemize}

\subsection{Оптимизированный алгоритм Винограда}

В оптимизированной версии применены две модификации:
\begin{enumerate}
    \item вынесение \texttt{size\_t half\_n = N / 2;} перед циклами. Затраты: $3$. Условия всех внутренних циклов принимают вид \texttt{k < half\_n}, ($f_{\text{усл}} = 1$), что исключает операцию деления на каждой итерации;
    \item вычисление вспомогательного индекса \texttt{size\_t two\_k = k << 1;} в начале каждой итерации внутреннего цикла ($1$ сдвиг + $1$ присваивание = $2$) и его повторное использование.
\end{enumerate}

Для векторов $row$ и $col$ внутреннее тело цикла: вычисление \texttt{two\_k} ($2$) и накопление \texttt{row[i] += A[i][two\_k] * A[i][two\_k + 1];}. Итого $f_{\text{тела}} = 2 + 9 = 11$.
Цикл по $k$ с условием ($f_{\text{усл}} = 1$): $f_k = 2 + H \cdot (2 + 11) = 13H + 2$.
Внешний цикл по строкам:
\begin{equation}
    f_{row\_\text{опт}} = 2 + M \cdot (2 + 13H + 2) = 13MH + 4M + 2.
\end{equation}
Аналогично для вектора $col$:
\begin{equation}
    f_{col\_\text{опт}} = 13QH + 4Q + 2.
\end{equation}

В основной части внутреннее тело цикла: вычисление \texttt{two\_k} ($2$) и выражение:
\begin{center}
\texttt{C[i][j] += (A[i][two\_k] + B[two\_k + 1][j]) * (A[i][two\_k + 1] + B[two\_k][j]);}
\end{center}
Итого $f_{\text{тела}} = 2 + (8 + 2 + 2 + 2 + 1) = 19$.

Цикл по $k$ с условием ($f_{\text{усл}} = 1$):
\begin{equation}
    f_{k\_\text{осн\_опт}} = 2 + H \cdot (2 + 19) = 21H + 2.
\end{equation}
С учётом $f_{\text{иниц\_}C} = 7$ цикл по $j$ даёт:
\begin{equation}
    f_{j\_\text{осн\_опт}} = 2 + Q \cdot (2 + 7 + 21H + 2) = 21QH + 11Q + 2.
\end{equation}
Внешний цикл по $i$:
\begin{equation}
    f_{\text{осн\_опт}} = 2 + M \cdot (2 + 21QH + 11Q + 2) = 21MQH + 11MQ + 4M + 2.
\end{equation}

Суммирование компонентов вместе с предварительным расчётом $f_{\text{кэш}} = 3$ и этапом нечётной коррекции $f_{\text{нечёт}}$:
\begin{equation}
    f_{\text{опт}} = 21MQH + 13H(M + Q) + 11MQ + 8M + 4Q + 9 + f_{\text{нечёт}}.
\end{equation}

Для квадратных матриц ($M = N = Q$):
\begin{itemize}
    \item при чётном $N$ ($H = N / 2$, $f_{\text{нечёт}} = 3$):
    \begin{equation}
        f_{\text{опт\_чёт}}(N) = 10{,}5N^3 + 24N^2 + 12N + 12;
    \end{equation}
    \item при нечётном $N$ ($H = (N - 1) / 2$, $f_{\text{нечёт}} = 13N^2 + 4N + 5$):
    \begin{equation}
        f_{\text{опт\_нечёт}}(N) = 10{,}5N^3 + 26{,}5N^2 + 3N + 14.
    \end{equation}
\end{itemize}

\subsection{Сводная аналитическая оценка}

В таблице~\ref{tbl:theoretical-complexity} приведены полученные аналитические оценки трудоёмкости для всех трёх исследуемых реализаций на квадратных матрицах $N \times N$.

\begin{table}[H]
    \caption{Теоретическая трудоёмкость алгоритмов умножения матриц}
    \label{tbl:theoretical-complexity}
    \centering
    \begin{tblr}{
            width = \linewidth,
            colspec = {X[1.6,l,m] X[2.3,c,m] X[2.3,c,m] Q[c,m]},
            hlines, vlines,
            row{1} = {font=\bfseries},
        }
        Алгоритм & Чётный размер $N$ & Нечётный размер $N$ & Сложность \\
        Стандартный & $11N^3 + 4N^2 + 4N + 2$ & $11N^3 + 4N^2 + 4N + 2$ & $O(N^3)$ \\
        Винограда (базовый) & $14{,}5N^3 + 30N^2 + 16N + 9$ & $14{,}5N^3 + 28{,}5N^2 + 3N + 11$ & $O(N^3)$ \\
        Винограда (оптим.) & $10{,}5N^3 + 24N^2 + 12N + 12$ & $10{,}5N^3 + 26{,}5N^2 + 3N + 14$ & $O(N^3)$ \\
    \end{tblr}
\end{table}

Согласно введённой модели вычислений, базовый алгоритм Винограда оказывается более трудоёмким по сравнению со стандартным алгоритмом ($14{,}5 \cdot N^3$ против $11 \cdot N^3$ при старшем члене): двукратное сокращение количества умножений элементов матрицы нивелируется многократным повторным умножением $2 \cdot k$ при индексации и повторным делением $N / 2$ в условиях циклов. Применение оптимизаций варианта (побитовый сдвиг и кэширование $N / 2$) снижает коэффициент при кубическом члене с $14{,}5$ до $10{,}5$, делая оптимизированный алгоритм Винограда наименее трудоёмким из всех трёх версий.

\section{Характеристики устройства}

Экспериментальные замеры процессорного времени выполнялись на вычислительном комплексе со следующими характеристиками:
\begin{itemize}
    \item процессор: Intel Core i5-1335U, 10 ядер, 12 потоков, тактовая частота до 4{,}60~ГГц;
    \item оперативная память: 16~ГБ DDR4, частота 3200~МГц;
    \item операционная система: Linux Ubuntu 24.04.4 LTS;
    \item компилятор: GCC 13.3.0, флаги сборки: \texttt{-std=c++20 -Wall -Wextra -Wpedantic -O0}.
\end{itemize}

\section{Инструменты замера процессорного времени}

Для измерения процессорного времени работы алгоритмов использовалась функция \texttt{clock()} из заголовочного файла \texttt{<ctime>}~\cite{cpp_clock}.

Перевод полученного значения в секунды осуществляется делением на константу \texttt{CLOCKS\_PER\_SEC}:
\begin{equation}
    t = \frac{\text{clock}_{\text{end}} - \text{clock}_{\text{start}}}{\text{CLOCKS\_PER\_SEC} \cdot K},
\end{equation}
где $K$ --- количество повторов умножения одной и той же пары матриц на текущем размере.

\section{Результаты массированных замеров}

Массированные замеры проводились для квадратных матриц в диапазоне размеров $N \in [10, 500]$ раздельно для чётных и нечётных размерностей. Результаты замеров сведены в таблицы~\ref{tbl:bench-even} и~\ref{tbl:bench-odd}.

\begin{table}[H]
    \caption{Процессорное время умножения матриц чётных размеров, с}
    \label{tbl:bench-even}
    \centering
    \begin{tblr}{
            width = \linewidth,
            colspec = {X[0.8,c,m] X[2,c,m] X[2,c,m] X[2,c,m]},
            hlines, vlines,
            row{1} = {font=\bfseries},
        }
        Размер $N$ & Стандартный & Виноград (базовый) & Виноград (оптим.) \\
        10  & 0{,}000029 & 0{,}000028 & 0{,}000028 \\
        20  & 0{,}000195 & 0{,}000119 & 0{,}000136 \\
        40  & 0{,}000501 & 0{,}000430 & 0{,}000427 \\
        60  & 0{,}001311 & 0{,}001398 & 0{,}001395 \\
        80  & 0{,}004431 & 0{,}003607 & 0{,}003615 \\
        100 & 0{,}005944 & 0{,}006139 & 0{,}006177 \\
        150 & 0{,}019987 & 0{,}020349 & 0{,}020742 \\
        200 & 0{,}047282 & 0{,}047742 & 0{,}048545 \\
        250 & 0{,}092496 & 0{,}093359 & 0{,}094982 \\
        300 & 0{,}161527 & 0{,}159662 & 0{,}163999 \\
        400 & 0{,}377031 & 0{,}374348 & 0{,}394725 \\
        500 & 0{,}734502 & 0{,}756302 & 0{,}766590 \\
    \end{tblr}
\end{table}

\begin{table}[H]
    \caption{Процессорное время умножения матриц нечётных размеров, с}
    \label{tbl:bench-odd}
    \centering
    \begin{tblr}{
            width = \linewidth,
            colspec = {X[0.8,c,m] X[2,c,m] X[2,c,m] X[2,c,m]},
            hlines, vlines,
            row{1} = {font=\bfseries},
        }
        Размер $N$ & Стандартный & Виноград (базовый) & Виноград (оптим.) \\
        11  & 0{,}000010 & 0{,}000011 & 0{,}000012 \\
        21  & 0{,}000061 & 0{,}000083 & 0{,}000074 \\
        41  & 0{,}000423 & 0{,}000462 & 0{,}000472 \\
        61  & 0{,}001371 & 0{,}001473 & 0{,}001461 \\
        81  & 0{,}003182 & 0{,}003370 & 0{,}003348 \\
        101 & 0{,}006288 & 0{,}006358 & 0{,}006437 \\
        151 & 0{,}020745 & 0{,}020833 & 0{,}021056 \\
        201 & 0{,}047889 & 0{,}048252 & 0{,}050058 \\
        251 & 0{,}093132 & 0{,}093312 & 0{,}101094 \\
        301 & 0{,}164249 & 0{,}174337 & 0{,}168708 \\
        401 & 0{,}381010 & 0{,}395288 & 0{,}393645 \\
        501 & 0{,}741757 & 0{,}764204 & 0{,}772992 \\
    \end{tblr}
\end{table}

Графики зависимости процессорного времени от линейного размера квадратных матриц для чётных и нечётных размеров приведены на рисунках~\ref{fig:plot_even} и~\ref{fig:plot_odd}.

\begin{figure}[H]
    \centering
    \begin{tikzpicture}
        \begin{axis}[
            width=0.92\linewidth,
            height=0.55\columnwidth,
            xlabel={Линейный размер матрицы $N$},
            ylabel={Процессорное время, с},
            grid=both,
            grid style={line width=.1pt, draw=gray!20},
            major grid style={line width=.2pt,draw=gray!50},
            legend pos=north west,
            legend cell align={left}
        ]
            \addplot[blue, mark=*, thick] coordinates {
                (10, 0.000029) (20, 0.000195) (40, 0.000501) (60, 0.001311) 
                (80, 0.004431) (100, 0.005944) (150, 0.019987) (200, 0.047282) 
                (250, 0.092496) (300, 0.161527) (400, 0.377031) (500, 0.734502)
            };
            \addlegendentry{Стандартный алгоритм}

            \addplot[red, mark=square*, thick] coordinates {
                (10, 0.000028) (20, 0.000119) (40, 0.000430) (60, 0.001398) 
                (80, 0.003607) (100, 0.006139) (150, 0.020349) (200, 0.047742) 
                (250, 0.093359) (300, 0.159662) (400, 0.374348) (500, 0.756302)
            };
            \addlegendentry{Базовый Виноград}

            \addplot[green!60!black, mark=triangle*, thick] coordinates {
                (10, 0.000028) (20, 0.000136) (40, 0.000427) (60, 0.001395) 
                (80, 0.003615) (100, 0.006177) (150, 0.020742) (200, 0.048545) 
                (250, 0.094982) (300, 0.163999) (400, 0.394725) (500, 0.766590)
            };
            \addlegendentry{Оптимизированный Виноград}
        \end{axis}
    \end{tikzpicture}
    \caption{Зависимость процессорного времени от размера матрицы $N$ (чётные размеры)}
    \label{fig:plot_even}
\end{figure}

\begin{figure}[H]
    \centering
    \begin{tikzpicture}
        \begin{axis}[
            width=0.92\linewidth,
            height=0.55\columnwidth,
            xlabel={Линейный размер матрицы $N$},
            ylabel={Процессорное время, с},
            grid=both,
            grid style={line width=.1pt, draw=gray!20},
            major grid style={line width=.2pt,draw=gray!50},
            legend pos=north west,
            legend cell align={left}
        ]
            \addplot[blue, mark=*, thick] coordinates {
                (11, 0.000010) (21, 0.000061) (41, 0.000423) (61, 0.001371) 
                (81, 0.003182) (101, 0.006288) (151, 0.020745) (201, 0.047889) 
                (251, 0.093132) (301, 0.164249) (401, 0.381010) (501, 0.741757)
            };
            \addlegendentry{Стандартный алгоритм}

            \addplot[red, mark=square*, thick] coordinates {
                (11, 0.000011) (21, 0.000083) (41, 0.000462) (61, 0.001473) 
                (81, 0.003370) (101, 0.006358) (151, 0.020833) (201, 0.048252) 
                (251, 0.093312) (301, 0.174337) (401, 0.395288) (501, 0.764204)
            };
            \addlegendentry{Базовый Виноград}

            \addplot[green!60!black, mark=triangle*, thick] coordinates {
                (11, 0.000012) (21, 0.000074) (41, 0.000472) (61, 0.001461) 
                (81, 0.003348) (101, 0.006437) (151, 0.021056) (201, 0.050058) 
                (251, 0.101094) (301, 0.168708) (401, 0.393645) (501, 0.772992)
            };
            \addlegendentry{Оптимизированный Виноград}
        \end{axis}
    \end{tikzpicture}
    \caption{Зависимость процессорного времени от размера матрицы $N$ (нечётные размеры)}
    \label{fig:plot_odd}
\end{figure}

\section{Сравнительный анализ}

Для сопоставления результатов было проведено сравнение относительной трудоёмкости по теоретической модели $f_{\text{отн}}$ и замеренному процессорному времени $t_{\text{отн}}$ относительно стандартного алгоритма:
\begin{equation}
    f_{\text{отн}} = \frac{f_X(N)}{f_{\text{станд}}(N)}, \quad t_{\text{отн}} = \frac{t_X(N)}{t_{\text{станд}}(N)}.
\end{equation}

В таблице~\ref{tbl:rel-comparison} сведены значения для чётного размера $N = 500$.

\begin{table}[H]
    \caption{Сравнение алгоритмов со стандартным при $N = 500$}
    \label{tbl:rel-comparison}
    \centering
    \begin{tblr}{
            width = \linewidth,
            colspec = {X[2,l,m] X[2,c,m] X[1,c,m] X[2,c,m] X[1,c,m]},
            hlines, vlines,
            row{1} = {font=\bfseries},
        }
        Алгоритм & Трудоёмкость & $f_{\text{отн}}$ & Время, с & $t_{\text{отн}}$ \\
        Стандартный & $1\,376\,002\,002$ & $1{,}000$ & $0{,}7345$ & $1{,}000$ \\
        Виноград (базовый) & $1\,820\,008\,009$ & $1{,}323$ & $0{,}7563$ & $1{,}030$ \\
        Виноград (оптим.) & $1\,318\,506\,012$ & $0{,}958$ & $0{,}7666$ & $1{,}044$ \\
    \end{tblr}
\end{table}

Анализ теоретических и экспериментальных данных позволяет сделать следующие выводы:
\begin{itemize}
    \item При увеличении размера матриц в 2 раза (например, со $100$ до $200$ и с $250$ до $500$) время работы всех реализаций возрастает примерно в 8 раз ($0{,}0925$~с $\to$ $0{,}7345$~с, то есть в $7{,}94$ раза), что подтверждает кубический порядок роста $O(N^3)$ ($2^3 = 8$).
    
    \item Согласно теоретической модели вычислений, базовый алгоритм Винограда более трудоёмок, чем стандартный ($f_{\text{отн}} \approx 1{,}32$), что подтверждается экспериментально: стандартный алгоритм выполняется быстрее базовой версии Винограда ($0{,}7345$~с против $0{,}7563$~с).
    
    \item Разница по времени между всеми тремя реализациями при $N = 500$ составляет не более $4{,}4\%$ ($0{,}7345$~с, $0{,}7563$~с и $0{,}7666$~с), а оптимизированная версия показывает незначительное отставание от базовой.

    \item Качественных различий между чётными и нечётными размерами матриц не обнаружено. Поправка на нечётность $N$ имеет порядок $O(N^2)$ и при больших размерах ($N \ge 100$) вносит пренебрежимо малый вклад на фоне кубической составляющей $O(N^3)$.
\end{itemize}

\section{Вывод}

В исследовательской части была разработана модель вычислений, выполнен теоретический расчёт трудоёмкости всех исследуемых алгоритмов, проведены экспериментальные замеры процессорного времени и выполнен сравнительный анализ алгоритмов умножения матриц.